\begin{abstract}
Recognizing cricket batting shots automatically from broadcast television videos is a challenging computer vision and image processing problem because the shot label depends on how the batter's body and bat move rather than static appearance alone. Broadcast clips contain moving spectators, television scoreboards, harsh stadium shadows, rapid camera panning, multiple players on the pitch (bowler, umpire, wicketkeeper, and non-striker), and variations in batter handedness. In this project, we systematically explore and compare three progressive approaches to solve this problem:
(1)~\textbf{Approach~1 (Classical Image Processing + Cross-Modal Fusion)}, which uses CIE-LAB CLAHE illumination normalization, HSV color-space turf/pitch segmentation for crowd suppression, Bobick--Davis Motion History Images (MHI), camera-pan-compensated Farneb\"{a}ck optical flow (36-bin polar HOOF), and Canny--Hough bat angle histograms fused with a lightweight CNN via Gated Cross-Modal Attention;
(2)~\textbf{Approach~2 (End-to-End Spatio-Temporal ResNet50 + LSTM)}, which fine-tunes higher-level ResNet50 visual features across 16 uniformly sampled RGB frames with a recurrent LSTM sequence classifier; and
(3)~\textbf{Approach~3 (Pose-Guided Striker Isolation and Biomechanical Dual-GRU Fusion)}, which builds on the CricShotNet pipeline by tracking the striker across every frame via a Viterbi dynamic programming search over YOLO11m-Pose and bat detections, extracting a 68-dimensional torso-normalized skeleton motion descriptor, and fusing a skeleton GRU-128 with a frozen EfficientNetV2-S RGB GRU-128 under handedness-invariant mirror augmentation.
Experiments on the 10-class \textbf{CricShot10} dataset (1,894 clips) and the 15-class \textbf{CricShot10k} benchmark (10,091 clips across 3 leak-free match-level splits) show that Approach~1 gives the fastest CPU inference ($\sim$1.28\,s/clip) and 93\%--100\% recall on geometry-dominated strokes (\textit{Defense}, \textit{Lofted}, \textit{Straight Drive}, \textit{Sweep}), Approach~2 achieves 79.0\% test accuracy on full-frame 10-class clips, and in Approach~3 the 68-D skeleton stream raises Macro-F1 from \textbf{67.18\%} (frozen RGB baseline) to \textbf{75.93\%} alone (\textbf{+8.75 points}) and \textbf{77.90\%} (\textbf{78.72\% accuracy}) when fused with mirroring (\textbf{+10.72 points} paired gain). We also report honest negative findings on wrist-speed impact phase alignment and \textit{Pull} vs.\ \textit{Hook} confusion.
\end{abstract}
